\documentclass[aps,prc,twocolumn,showpacs,superscriptaddress,floatfix]{revtex4-2}
\usepackage{amsmath,amssymb,graphicx,bm}
\usepackage[colorlinks=true,linkcolor=blue,citecolor=blue]{hyperref}

\newcommand{\Esym}{E_{\mathrm{sym}}}
\newcommand{\Lsym}{L_{\mathrm{sym}}}
\newcommand{\DRnp}{\Delta R_{np}}
\newcommand{\Msun}{M_\odot}

\begin{document}

\title{Neutron skins of $^{208}$Pb and $^{48}$Ca predicted from
cooling-constrained relativistic mean field models}

\author{Nisha Singh}
\affiliation{Department of Physics, Shiv Nadar Institution of Eminence,
Gautam Buddha Nagar, Uttar Pradesh 201314, India}

\author{Rana Nandi}
\affiliation{Department of Physics, Indian Institute of Technology Delhi,
Hauz Khas, New Delhi 110016, India}

\date{\today}

\begin{abstract}
We compute the neutron skin thicknesses of $^{208}$Pb and $^{48}$Ca directly
from the relativistic mean field parametrizations whose equations of state are
constrained by neutron star thermal evolution, using the same coupling
constants throughout. The finite-nucleus problem is solved in spherical
Hartree approximation with no adjustment of parameters. Validated against
published results for three modern parametrizations at their own published
$(\Esym,\Lsym)$, the solver reproduces neutron skins to $0.7$--$4.5$
milli-fermi, binding energies to $0.08$~MeV per nucleon and charge radii to
$0.03$~fm, and returns the correct negative skins for the $N=Z$ nuclei. The
predicted $\DRnp$--$\Lsym$ slope for $^{208}$Pb passes through the published
model-averaged correlation slope within the sampled range, without having been
fitted to it.

Confronting the predictions with PREX-2 and CREX, we find that cooling does not
discriminate between them: the joint-$\chi^2$ cost mildly favours PREX-2
($0.9\sigma$ against $1.4\sigma$), while the marginal posteriors favour CREX,
whose $\Esym$ interval the cooling constraint overlaps and PREX-2's it does not.
We report both, and show the disagreement is a property of the model family. The underlying obstruction is structural:
the models predict nearly equal skins for the two nuclei, $0.177$ and
$0.174$~fm at the joint best fit, whereas the measurements call for $0.283$ and
$0.121$~fm. We also point out that the skin cannot be computed at all for
parametrizations fitted to uniform nuclear matter alone, because such fits
determine only the ratios $g_i/m_i$, leaving the meson masses -- and hence the
range of the scalar field and the diffuseness of the nuclear surface --
undefined.
\end{abstract}

\pacs{21.10.Gv, 21.65.Ef, 25.30.Bf, 26.60.Kp}
\keywords{neutron skin; symmetry energy; PREX; CREX; relativistic mean field}

\maketitle

% =====================================================================
\section{Introduction}

The neutron skin thickness
$\DRnp = \langle r^2\rangle_n^{1/2} - \langle r^2\rangle_p^{1/2}$
of a heavy nucleus is among the most direct laboratory probes of the symmetry
energy slope $\Lsym$ \cite{Brown2000,Furnstahl2002,RocaMaza2011}, and
parity-violating electron scattering \cite{HorowitzPiekarewicz2001} has now
measured it for two nuclei. PREX-2 reports
$\DRnp(^{208}\mathrm{Pb}) = 0.283\pm0.071$~fm \cite{PREX2}, while CREX reports
$\DRnp(^{48}\mathrm{Ca}) = 0.121\pm0.026\,(\mathrm{exp})\pm0.024\,(\mathrm{model})$~fm
\cite{CREX}.
Taken at face value the two results pull in opposite directions: a thick lead
skin favours a stiff symmetry energy \cite{Reed2021}, a thin calcium skin a
soft one, and no mean field family reproduces both \cite{Reinhard2022}.

Astrophysical constraints on $\Lsym$ are usually compared with these
measurements through a published $\DRnp$--$\Lsym$ correlation fitted across many
models \cite{RocaMaza2011}. That procedure carries the intrinsic scatter of the
fit, of order $\pm0.02$~fm, and more importantly it is not a prediction of the
particular Lagrangian whose equation of state was constrained. Here we remove
that step: for every parametrization used in the cooling analysis of the
companion paper \cite{PaperI} we solve the finite-nucleus problem with the
\emph{same} couplings, so the skin becomes an output of the same Lagrangian
rather than an input borrowed from elsewhere.

% =====================================================================
\section{Method}

\subsection{Field equations}

The Lagrangian is the standard non-linear relativistic mean field form
\cite{Walecka1974,BogutaBodmer1977,SerotWalecka1986} with
$\sigma$, $\omega$ and $\rho$ mesons, cubic and quartic $\sigma$
self-interactions, a quartic $\omega$ self-interaction $\zeta$
\cite{MullerSerot1996}, and an
$\omega$--$\rho$ cross coupling $\Lambda_{\omega\rho}$
\cite{HorowitzPiekarewicz2001}. In a finite nucleus the
meson fields acquire gradients and, writing $\Phi=g_\sigma\sigma$,
$W=g_\omega\omega_0$, $R=g_\rho\rho_{03}$, obey
\begin{align}
-\nabla^2\Phi + m_\sigma^2\Phi &= g_\sigma^2\big[n_s - bM\Phi^2 - c\Phi^3\big],\\
-\nabla^2 W + m_\omega^2 W &= g_\omega^2\Big[n_v - 2\Lambda_{\omega\rho}R^2W
   - \tfrac{\zeta}{6}W^3\Big],\\
-\nabla^2 R + m_\rho^2 R &= g_\rho^2\Big[\tfrac{n_3}{2}
   - 2\Lambda_{\omega\rho}W^2R\Big].
\end{align}
Setting $\nabla^2\to 0$ recovers the uniform-matter equations exactly, which is
the consistency check that these are the correct equations. Nucleons obey the
radial Dirac equations \cite{HorowitzSerot1981}
\begin{align}
\frac{dG}{dr} &= -\frac{\kappa}{r}G + \big[\varepsilon - V(r) + M^\ast(r)\big]F,\\
\frac{dF}{dr} &= +\frac{\kappa}{r}F - \big[\varepsilon - V(r) - M^\ast(r)\big]G,
\end{align}
with $M^\ast = M-\Phi$ and $V = W \pm R/2 + V_{\rm Coul}$. Densities follow from
the occupied orbitals in the usual way.

\subsection{Numerical treatment}

Single-particle energies are located by a Sturm-sequence bisection: integrating
the regular solution outward from the origin to the box edge, the number of
nodes of $G$ equals the number of eigenvalues below the trial energy. Because
that count is a monotone function of energy, plain bisection on a single integer
cannot skip or mislabel a level. We note that restricting the node count to the
nuclear interior, in an attempt to avoid the divergent tail, is incorrect: the
zero crossing in the tail is a genuine Sturm node, and discarding it shifts
every level by one.

The Klein--Gordon equations are solved with their exact Green's functions rather
than by finite differences, with the exponentials paired so that no
intermediate overflows. The nucleon basis is taken generously larger than the
occupied space; at the soft corner of the parameter grid the bound levels of
$^{208}$Pb hold exactly $126$ neutrons with no margin, and a basis chosen to
match the occupied orbitals fails there.

\subsection{The meson masses are not optional}
\label{sec:msigma}

A point of principle arises that has practical consequences. In uniform nuclear
matter only the ratios $C_i^2=(g_i/m_i)^2$ enter, so $m_\sigma$ and $g_\sigma$
are not separately determined. A finite nucleus is different: $m_\sigma$ fixes
the range $1/m_\sigma$ of the scalar field, which sets the diffuseness of the
nuclear surface and hence the radius.

Using a single generic $m_\sigma = 550$~MeV for parametrizations whose
published value is near $495$~MeV produces a surface that is too sharp. Across
the three parametrizations and both nuclei the charge radius falls short by
$0.06$--$0.12$~fm and the binding energy gains a spurious
$0.50$--$1.19$~MeV per nucleon, the larger errors being for $^{48}$Ca. We therefore adopt each parametrization's own meson masses as
tabulated in Ref.~\cite{Das2021}, holding
$C_\sigma^2$ fixed and rescaling $g_\sigma = C_\sigma m_\sigma$, which leaves
the uniform-matter limit -- and hence every equation of state and cooling
curve of Ref.~\cite{PaperI} -- unchanged while correcting the range.

\subsection{Parametrizations fitted to uniform matter cannot give a skin}
\label{sec:gm}

The same observation excludes an entire class of parametrization. Fits
performed against uniform nuclear matter alone publish only the ratios
$(g_i/m_i)^2$, because that is all uniform matter requires. The GM1 and GM2
sets of Ref.~\cite{GM1991} are of exactly this kind: Table~II of that paper
gives $(g/m)^2 = 11.79$, $7.149$ and $4.411$~fm$^2$ with $b=0.002947$ and
$c=-0.001070$, and no individual meson mass appears anywhere in it. For such
parametrizations $m_\sigma$ is genuinely undefined, and with it the range of the
scalar field and the diffuseness of the nuclear surface. A neutron skin
computed from them is a property of whatever meson mass one invents, not of the
model.

This is a statement about the completeness of the parametrization rather than a
defect: the fits are correct in the domain for which they were constructed. But
it means that two of the five parametrizations in our cooling analysis are
unusable here, and we exclude them. Independently, both also fail numerically
for $^{208}$Pb: their weak scalar fields, $M^\ast/M = 0.70$ and $0.78$, give a
weak spin--orbit force, the orbitals required to reach $N=126$ do not bind, and
the solution collapses.

To say anything more quantitative one has to assume a meson mass, which is the
very step this section argues against; we do so here only to characterise the
failure. With a generic $m_\sigma = 550$~MeV, GM1 converges for $^{16}$O,
$^{40}$Ca and $^{48}$Ca and is overbound by $0.52$--$0.62$~MeV per nucleon,
while GM2 converges only for $^{16}$O and $^{40}$Ca and is overbound by
$2.64$ and $1.23$~MeV per nucleon respectively. Those numbers are properties of
the assumed $550$~MeV, not of the parametrizations, which is precisely the
point: without a published meson mass there is no model-specific answer to
quote.

% =====================================================================
\section{Validation}

\begin{figure}[t]
\includegraphics[width=\columnwidth]{figs/fig_skin_validation.pdf}
\caption{Predicted neutron skins against published values \cite{Das2021} for the
three usable parametrizations, each evaluated at its own published
$(\Esym,\Lsym)$. Filled symbols $^{208}$Pb, open symbols $^{48}$Ca; the line is
exact agreement. No parameter is adjusted anywhere in the comparison.}
\label{fig:valid}
\end{figure}

\begin{table}[t]
\caption{Comparison with the published results of Ref.~\cite{Das2021} at each
parametrization's own published $(\Esym,\Lsym)$, taken from Table~2 of that
work. No parameter is adjusted. Published skins are from its Table~3.}
\label{tab:valid}
\begin{ruledtabular}
\begin{tabular}{lcccccc}
 & \multicolumn{2}{c}{}
 & \multicolumn{2}{c}{$\DRnp(^{208}$Pb$)$ (fm)}
 & \multicolumn{2}{c}{$\DRnp(^{48}$Ca$)$ (fm)}\\
Set & $\Esym$ & $\Lsym$ & this work & Ref.~\cite{Das2021}
    & this work & Ref.~\cite{Das2021}\\
\hline
BigApple  & 31.32 & 39.80 & 0.1535 & 0.151 & 0.1707 & 0.170\\
FSUGarnet & 30.95 & 51.04 & 0.1628 & 0.162 & 0.1655 & 0.169\\
IOPB-I    & 33.30 & 63.58 & 0.2226 & 0.221 & 0.1975 & 0.202\\
\end{tabular}
\end{ruledtabular}
\end{table}

Figure~\ref{fig:valid} and Table~\ref{tab:valid} compare predicted skins with
the published values of Ref.~\cite{Das2021}. Agreement
is $0.7$--$4.5$ milli-fermi, between fifteen and a hundred times finer than the
PREX-2 uncertainty. Charge radii agree to $0.03$~fm and binding energies to
$0.08$~MeV per nucleon, the latter a systematic under-binding consistent with
the pairing term included in the published calculation and omitted here; it does
not affect the skins, which are differences of radii.

Three further checks were not fitted to anything. First, the solver reproduces
one parametrization across $A=16$ to $208$ ($^{16}$O, $^{40}$Ca, $^{48}$Ca,
$^{208}$Pb) with no divergence and binding energies within
$0.09$~MeV per nucleon of experiment, the largest departure being $^{40}$Ca at
$-0.089$~MeV per nucleon. Second, the
$N=Z$ nuclei come out with \emph{negative} skins, $-0.030$~fm for $^{16}$O and
$-0.053$~fm for $^{40}$Ca, which is the correct sign because the Coulomb
repulsion pushes protons outward. Third, the local slope
$d\DRnp/d\Lsym$ for $^{208}$Pb falls from
$0.0024$--$0.0031$~fm/MeV at $\Lsym = 40$~MeV to
$0.0003$--$0.0007$~fm/MeV at $\Lsym = 100$~MeV as the skin saturates, passing
through the published model-averaged correlation slope of $0.00147$~fm/MeV
\cite{RocaMaza2011} near $\Lsym \simeq 80$~MeV; in the band the cooling
constraint selects, $\Lsym = 50$--$75$~MeV, it is $0.0017$--$0.0023$~fm/MeV. We
report this as consistency rather than as a reproduction of the published
correlation, because a within-model derivative and a cross-model regression
slope are not the same quantity; no fit to that correlation was performed at any
stage.

\subsection{Interpolation across the grid}

For the parameter scan the computed skins are represented by a surface
$\DRnp = a + b\Lsym + c\Lsym^2 + d\Esym + e\,\Esym\Lsym$. The quadratic term in
$\Lsym$ is necessary rather than cosmetic: $\DRnp$ rises steeply with $\Lsym$
and then saturates, as the slopes above show, and a purely linear representation
leaves root-mean-square residuals of $0.015$--$0.018$~fm for $^{208}$Pb,
half the PREX-2 uncertainty. With the quadratic term the residuals fall to
$0.0034$--$0.0073$~fm, which we add in quadrature to the
experimental errors throughout.

$\Esym = 25$~MeV is excluded from the scan: at that corner $^{208}$Pb does not
bind, and the near-threshold single-particle states there are in any case
outside the domain where a Hartree treatment without pairing is reliable.

\begin{figure}[t]
\includegraphics[width=\columnwidth]{figs/fig_skins.pdf}
\caption{Predicted $\DRnp$ against $\Lsym$ at $\Esym=32.5$~MeV, with the PREX-2
\cite{PREX2} and CREX \cite{CREX} central values and uncertainties shaded. The
models cannot simultaneously produce a thick $^{208}$Pb skin and a thin
$^{48}$Ca skin.}
\label{fig:skins}
\end{figure}

% =====================================================================
\section{Confrontation with PREX-2 and CREX}

\subsection{The tension is reproduced, not imported}

Figure~\ref{fig:skins} shows the predictions against both measurements.
At the point that best satisfies both simultaneously the
predictions are $\DRnp(^{208}\mathrm{Pb}) \simeq 0.177$~fm and
$\DRnp(^{48}\mathrm{Ca}) \simeq 0.174$~fm: nearly equal. The measurements call
for $0.283$ and $0.121$~fm. The joint $\chi^2$ at that point is $4.42$. With two
measurements and two free parameters this quantity has no degrees of freedom, so
it is not a goodness of fit; we use it only as a \emph{relative} measure of how
far apart the two measurements pull within this model family, and we deliberately
refrain from converting it into a significance, which would require a
degrees-of-freedom count that does not exist here.

This is the PREX--CREX problem \cite{Reinhard2022} appearing within our own
calculation rather than being taken over from elsewhere. Mean field models tie
the two skins together through the same isovector couplings
\cite{Furnstahl2002}, and this family is no exception.

\subsection{Constraints from each measurement separately}

Comparing our predicted skins with each measured skin, without reference to any
published $\Lsym$ value, gives at 68\% credibility
\begin{align}
\text{PREX-2 alone:}\quad & \Lsym = 50\text{--}95~\mathrm{MeV},\
\Esym = 35.5\text{--}43~\mathrm{MeV},\\
\text{CREX alone:}\quad & \Lsym = 50\text{--}70~\mathrm{MeV},\
\Esym = 26.5\text{--}32.5~\mathrm{MeV}.
\end{align}
The two do not overlap in $\Esym$. Our PREX-2-alone $\Lsym$ range is consistent
with, and narrower than, the $106\pm37$~MeV of Ref.~\cite{Reed2021}, the
difference being that our $\Esym$ is free to move rather than held at a
calibrated value.

\begin{table}[t]
\caption{Constraints on $(\Esym,\Lsym)$ at 68\% credibility, all derived within
the same three parametrizations and over the same reduced grid. The skin
constraints use only our predicted skins and the measured skins; no published
$\Lsym$ value from either experiment enters at any point.}
\label{tab:constraints}
\begin{ruledtabular}
\begin{tabular}{lcc}
Probe & $\Lsym$ (MeV) & $\Esym$ (MeV)\\
\hline
cooling alone\footnotemark[1] & $[45,\,70]$ & $[26.5,\,34.0]$\\
PREX-2 alone             & $[50,\,95]$ & $[35.5,\,43.0]$\\
CREX alone               & $[50,\,70]$ & $[26.5,\,32.5]$\\
cooling $+$ PREX-2       & $[50,\,75]$ & $[28.0,\,35.5]$\\
cooling $+$ CREX         & $[50,\,70]$ & $[26.5,\,34.0]$\\
cooling $+$ both\footnotemark[2] & $[45,\,70]$ & $[26.5,\,34.0]$\\
\end{tabular}
\end{ruledtabular}
\footnotetext[1]{This differs by one grid step from the
$\Lsym = 50$--$75$~MeV quoted in Ref.~\cite{PaperI}, which is not an
inconsistency but a consequence of the reduced grid used here: see
Sec.~\ref{sec:reconcile}.}
\footnotetext[2]{Shown for completeness only, and \emph{not} recommended as a
combined constraint. PREX-2 and CREX are mutually inconsistent within this
model family, so merging them into one posterior yields an interval that
describes neither measurement. That it comes out equal to the cooling-only row
is precisely the point: the two skin constraints pull in opposite directions
and very nearly cancel, leaving cooling to set the answer. The row should be
read as a demonstration of that cancellation, not as a result.}
\end{table}

\subsection{Reconciling the cooling interval with the companion paper}
\label{sec:reconcile}

The cooling-alone row of Table~\ref{tab:constraints} reads
$\Lsym = 45$--$70$~MeV, whereas Ref.~\cite{PaperI} quotes $50$--$75$~MeV from
the same cooling calculation. The difference is entirely one of grid, and we
state it explicitly so that it is not mistaken for a disagreement.

Reference~\cite{PaperI} marginalises over all five parametrizations and the
full $\Esym$ range. The present paper must drop GM1 and GM2, which have no
published meson masses (Sec.~\ref{sec:gm}), and must also drop
$\Esym = 25$~MeV, where $^{208}$Pb does not bind and no skin exists to
interpolate. Applying either restriction on its own leaves the interval at
$[50,75]$~MeV; it is only the two together, which remove $397$ of the $864$
equations of state and in particular the entire soft-$\Esym$ edge of the three
remaining parametrizations, that move it down by one $5$~MeV step. The two
intervals overlap over four fifths of their length, and the quantity being
constrained is the same.

For every statement in this paper that combines cooling with a skin
measurement, the cooling likelihood is evaluated on the reduced grid, so the
comparison is internally consistent.

\subsection{Combination with the cooling constraint}

Table~\ref{tab:constraints} collects the constraints from each probe.
Combining with the thermal constraint of Ref.~\cite{PaperI}, evaluated over the
same three parametrizations for consistency, and measuring the
tension as the $\chi^2$ cost of satisfying cooling and one experiment together,
beyond satisfying each separately, we obtain $0.78$ ($0.9\sigma$) for PREX-2 and
$2.07$ ($1.4\sigma$) for CREX.

By that measure cooling sits closer to PREX-2. The difference between the two
costs, however, is $1.29$ in $\chi^2$, or about $1.1\sigma$: it does not
discriminate.

\textbf{More importantly, the second natural measure points the other way, and
we report both.} Comparing marginal posteriors rather than joint cost, the
cooling $\Esym$ interval $[26.5,34.0]$~MeV does not overlap the PREX-2 interval
$[35.5,43.0]$ at all, while it overlaps almost the whole of the CREX interval
$[26.5,32.5]$. Adding CREX to the cooling constraint leaves the inferred
interval unchanged, whereas adding PREX-2 shifts it upward to $[28.0,35.5]$.
The peaks say the same thing: cooling peaks at $\Esym \simeq 32.5$~MeV against
$26.5$ for CREX and $41.5$ for PREX-2, so the cooling peak lies nearer CREX.

The two measures disagree for a reason worth stating, because it is a property
of the model family rather than of the data. PREX-2 can be reproduced
\emph{exactly} somewhere in the grid --- its minimum $\chi^2$ alone is $0.00$
--- but only at $\Esym \simeq 41.5$~MeV, a region cooling disfavours; the cost
of $0.78$ measures the degradation incurred in going there. CREX can never be
reproduced exactly: its minimum $\chi^2$ alone is $1.03$ anywhere on the grid.
So CREX begins with a penalty the model cannot remove, even though its preferred
region coincides with cooling's.

\textbf{Our conclusion is therefore that cooling does not discriminate between
the two measurements}, and that the apparent direction of any preference depends
on which statistic is chosen. We regard reporting a single one of these numbers
as the more serious error, and note that an earlier version of this work quoted
only the joint-cost comparison.

\subsection{Predicted skins at the cooling optimum}

\begin{table}[t]
\caption{Skins predicted at the joint optimum, compared with measurement. The
model cannot produce a thick $^{208}$Pb skin together with a thin $^{48}$Ca
skin; it produces nearly equal skins instead. The CREX uncertainty is the
experimental and model errors of Ref.~\cite{CREX} combined in quadrature.}
\label{tab:pred}
\begin{ruledtabular}
\begin{tabular}{lccc}
Nucleus & predicted (fm) & measured (fm) & difference\\
\hline
$^{208}$Pb & 0.177 & $0.283\pm0.071$ \cite{PREX2} & $-1.5\sigma$\\
$^{48}$Ca  & 0.174 & $0.121\pm0.035$ \cite{CREX}  & $+1.5\sigma$\\
\end{tabular}
\end{ruledtabular}
\end{table}

Table~\ref{tab:pred} states the central difficulty compactly. The predictions
sit almost exactly midway between the two measurements, missing each by a
comparable amount and in opposite directions. No adjustment of $\Esym$ or
$\Lsym$ within this family resolves that, because both skins respond to the same
isovector couplings in the same direction.

% =====================================================================
\section{Conclusions}

Solving the finite-nucleus problem with the same couplings that generate the
equation of state turns the neutron skin into a prediction of the constrained
Lagrangian rather than a quantity inferred through a model-averaged correlation.
The solver reproduces the published skins of Ref.~\cite{Das2021} to a few
milli-fermi, recovers the known $\DRnp$--$\Lsym$ behaviour without being fitted
to it, and returns the correct sign for $N=Z$ nuclei.

Applied to parametrizations constrained by neutron star cooling \cite{PaperI},
the predicted skins lie between the PREX-2 \cite{PREX2} and CREX \cite{CREX}
central values and exclude neither. Which experiment appears closer depends on
the statistic: the joint-$\chi^2$ cost favours PREX-2 at $0.9\sigma$ against
$1.4\sigma$, while the marginal $\Esym$ posteriors favour CREX. Cooling does
not discriminate. The
obstruction to doing better is structural rather than statistical: mean field
models predict nearly equal skins for $^{208}$Pb and $^{48}$Ca, and the two
measurements do not \cite{Reinhard2022}.

Finally, the requirement that a parametrization specify its meson masses, and
not merely the ratios $g_i/m_i$, is worth stating explicitly. Uniform nuclear
matter never needs them; a finite nucleus cannot be computed without them. Any
astrophysically constrained equation of state intended for comparison with
neutron skin measurements must therefore come from a parametrization complete in
this sense.

\begin{acknowledgments}
N.S. thanks Kenji Nishiwaki for supervision and support, and acknowledges Shiv
Nadar Institution of Eminence for a research fellowship. The numerical
calculations were performed on institutional computing resources at Shiv Nadar
Institution of Eminence.
\end{acknowledgments}

\bibliographystyle{apsrev4-2}
\bibliography{refs}

\end{document}
